\documentclass{article}

\usepackage{graphicx}
\usepackage[a4paper, margin=1in]{geometry}
\usepackage{amsmath}
\usepackage{IEEEtrantools}
\usepackage{karnaugh-map}
\usepackage{minted}

\graphicspath{{./assets/}}

\title{Homework Five}
\author{Aden Perry}
\date{\today}

\begin{document}
\maketitle

\begin{enumerate}
	\item The MDIS bit of the MCR register is set to safely (without faults)
	      configure the PIT. The PTI bit of SIM\_SCGC6 register is set to allow
	      the system clock to be used with the PIT module. The TIE bit of the
	      TCTRL register is cleared to allow configuration of interrupts. IRQ
	      22 of NVIC IPR 5 is set with the desired priority, 0--3. NVIC\_IPCR
	      is cleared to clear pending interrupts. The PIT bits of the
	      NVIC\_ISER are set to enable the PIT interrupt. MDIS is re-enabled to
	      enable module; configuration of load time is done; finally, the
	      transmit interrupt is enabled (TIE bit).

	\item The processor pushes R0-R3 and R12-15 to the stack. A vector fetch is
	      performed to determine the address to run the exception at.
	      Simultanously, an EXC\_RETURN value is written to LR. Control is
	      returned to the program when one of \texttt{0xFFFFFFF9} or
	      \texttt{0xFFFFFFFD} (\texttt{0xFFFFF1} returns to the previous
	      interrupt, not the program).

	\item They are values between 0 and 3 assigned to each interrupt and are
	      used in the event that two interrupts (or exceptions) occur
	      simultanously; the higher priority interrupt is dealt with first,
	      forcing late-arrival if necessary, and then tail-chains into the
	      lower priority interrupt.

	\item \inputminted{asm}{q4.s}

	\item The TPM0 bit of the SIM\_SCGC6 register is set to allow the system
	      clock to reach the TPM0 module. The PORTB bit of the SIM\_SCGC5 is
	      set to allow the system clock to reach Port B of the KL05. The MUX
	      field of the PORTB\_PCR register is set to 111 to select TMP0 channel
	      2.

	\item Negative, so set sign bit. 16.75 = 10000.11. Normalize: 1.000011. Put
	      into IEEE-754 and convert:
	      -16.75 = \texttt{1  1000 0011  0000 1100 0000 00000000 0000 0000 000} \\
	      -16.75 = \texttt{C186 0000}

	\item \texttt{0x40b40000} = \texttt{0 1000 0001 0110 1000 0000 0000 0000 000} \\
	      Sign bit is 0, number is $101.101_2$ = 5.625

	\item
	      \begin{enumerate}
		      \item \texttt{0xc0a8 0000} $\cdot$ \texttt{0x3fd0 0000} = \texttt{0xc108 8000}
		            \begin{align*}
			            c0a80000 \cdot 3fd00000 = &                                                                                        \\
			                                      & \left(1\right)\_\left(10000001\right)\_\left(01010000000000000000000\right)            \\
			                                      & \cdot \left(0\right)\_\left(01111111\right)\_\left(10100000000000000000000\right)      \\
			            %
			            =                         & \left(0 \text{ XOR } 1\right)                                                          \\
			                                      & \_\left(10000001+01111111-01111111\right)                                              \\
			                                      & \_\left(1.01010000000000000000000 \cdot 1.10100000000000000000000\right)               \\
			            %
			            =                         & \left(1\right) \_ \left(10000001\right)                                                \\
			                                      & \_\left(1.01010000000000000000000 \cdot 1.10100000000000000000000\right)               \\
			            %
			            =                         & \left(1\right) \_ \left(10000001\right)                                                \\
			                                      & \_\left(0.0010101 + 0.10101 + 1.0101\right)                                            \\
			            %
			            c108 8000=                & \left(1\right) \_ \left(100 0001 0\right) \_ \left(000 1000 1000 0000 0000 0000\right)
		            \end{align*}

		      \item \texttt{0xc0a8 0000} $+$ \texttt{0x3fd0 0000} =
		            \begin{align*}
			            c0a8 0000 + 3fd0 0000 & =                                                                                    \\
			                                  & \left(1\right)\_\left(100 0000 1\right)\_\left(010 1000 0000 0000 0000 0000\right)   \\
			                                  & + \left(0\right)\_\left(011 1111 1\right)\_\left(101 0000 0000 0000 0000 0000\right) \\
		            \end{align*}
	      \end{enumerate}

	\item
	      Can be done by scanning through the entire A6.7 section of the
	      ARM-v6M manual, because it can't be searched by the binary code, or
	      by:
	      \begin{verbatim}
            printf '\x23\x06\x46\x18\x81\x02\x10\x91' > a.out && \
            arm-none-eabi-objdump -EL -D -b binary -m arm -M force-thumb a.out
          \end{verbatim}
	      Either way, so long as little-endianess is considered, the following
	      result is achieved:

	      \begin{tabular}{ccc}
		      Hex  & Binary              & ASM                 \\
		      \hline
		      \hline
		      2306 & 0010 0011 0000 0110 & lsls r3, r4, \#24   \\
		      4618 & 0100 0110 0001 1000 & adds r6, r0, r1     \\
		      8102 & 1000 0001 0000 0010 & lsls r1, r0, \#10   \\
		      1091 & 0001 0000 1001 0001 & str  r1, [sp, \#64] \\
		      \hline
		      \hline
	      \end{tabular}

	\item
	      Observing little-endianess:

	      \begin{tabular}{cc}
		      ASM               & Hex  \\
		      \hline
		      \hline
		      movs r4, \#15     & 0f24 \\
		      mov r7, sp        & 6f46 \\
		      subs r0, r4, r6   & a01b \\
		      adds r2, r2, \#12 & 0c32 \\
		      \hline
		      \hline
	      \end{tabular}

	\item

	      \begin{enumerate}
		      \item
		            \begin{itemize}
			            \item ALU
			            \item Control Unit
			            \item CPU clock (from system clock)
			            \item L1-3 caches
			            \item Register file
		            \end{itemize}
		      \item
		            \begin{itemize}
			            \item Data bus
			            \item Control bus
			            \item Address bus
			            \item Communicates with CPU, RAM, and I/O. Per my
			                  understanding, they're connected to but not part
			                  of the bus.
		            \end{itemize}
		      \item
		            \begin{enumerate}
			            \item Registers
			            \item L1 Cache
			            \item L2 Cache
			            \item L3 Cache
			            \item Main memory
			            \item HDD/SDD
			            \item DVD ROM
		            \end{enumerate}
	      \end{enumerate}
\end{enumerate}

\end{document}
